\section{Conclusions}
\label{sec:conclusions}

We study when multilingual evaluation supports decisions during pretraining. Predictable learning curves, evidence above chance, and agreement with a larger model answer different questions. HellaSwag, LAMBADA, and per-language BPB track scale consistently, while reformulating letter-based tasks as answer continuations reveals signal that their original protocols suppress. Across the 12-checkpoint evaluation grid, BPB rankings transfer more consistently to the 1.7B reference than the aggregate benchmark rankings, but early agreement can reverse during training. Cross-task transfer is limited, and the most informative SNR definition varies by language and by whether the decision concerns model size or training progress. These findings favor task- and language-specific checks over a single benchmark score or universal proxy threshold.

Our next steps address two questions left open by this analysis. First, we will test how predictivity depends on the design intervention, comparing language selection, sampling temperature, second-language choice, and architecture while estimating uncertainty in the reference ranking from repeated runs. Second, we will test whether scaling behavior and decision reliability transfer to held-out languages, with explicit separation between languages absent from pretraining and languages withheld from fitting the predictor. These studies will also provide held-out decisions for calibrating surrogate-based selection and stopping rules, extending the present descriptive evidence toward actionable guidance for multilingual model development.
